% Original analytic teaching diagram, 2026-11-08.
\definecolor{VB}{HTML}{7188A3}\definecolor{VBF}{HTML}{E5EAF0}
\definecolor{VG}{HTML}{829C8B}\definecolor{VGF}{HTML}{E8EEE9}
\definecolor{VY}{HTML}{B4A16F}\definecolor{VYF}{HTML}{F1ECDD}
\definecolor{VR}{HTML}{A7787D}\definecolor{VRF}{HTML}{F0E5E6}
\definecolor{VP}{HTML}{9588A5}\definecolor{VPF}{HTML}{ECE8F0}
\definecolor{VI}{HTML}{3B4252}\definecolor{VM}{HTML}{6B7280}\definecolor{VL}{HTML}{4C4D4F}
\begin{tikzpicture}[x=1mm,y=1mm,>=Stealth,font=\figurefont\fontsize{8}{10}\selectfont,text=VI,
lab/.style={text=VM,font=\figurefont\fontsize{7.5}{9}\selectfont},
pt/.style={circle,fill=VI,inner sep=0pt,minimum size=1.5mm}]
% unit 20mm per length unit; side .9 ; circumradius .9/sqrt 3
\def\S{18}\def\Rc{10.392}\def\H{15.588}
\foreach \cx/\tag in {20/(a),58/(b),96/(c)} {
  \coordinate (A\cx) at (\cx-9,18);
  \coordinate (B\cx) at (\cx+9,18);
  \coordinate (C\cx) at (\cx,18+\H);
  \node[anchor=base] at (\cx,-2.4) {\tag};
}
% (a) balls of radius .5 (10mm); the centre lies at distance .520 and escapes
\foreach \p in {A20,B20,C20} { \fill[VBF] (\p) circle (10); }
\foreach \p in {A20,B20,C20} { \draw[VB,line width=.6pt] (\p) circle (10); }
\draw[VY,line width=.8pt,dash pattern=on 2pt off 1.6pt] (20,18+5.196) circle (\Rc);
\node[circle,fill=VR,inner sep=0pt,minimum size=1.3mm] at (20,18+5.196) {};
\foreach \p in {A20,B20,C20} { \node[pt] at (\p) {}; }
\node[lab,anchor=north] at (20,6.6) {$r=0.5$};
% (b) Cech
\draw[VG,line width=1.2pt] (A58)--(B58)--(C58)--cycle;
\foreach \p in {A58,B58,C58} { \node[pt] at (\p) {}; }
\node[lab,anchor=north] at (58,6.6) {$\check C_{0.5}$};
% (c) Rips
\fill[VYF] (A96)--(B96)--(C96)--cycle;
\draw[VG,line width=1.2pt] (A96)--(B96)--(C96)--cycle;
\foreach \p in {A96,B96,C96} { \node[pt] at (\p) {}; }
\node[lab,anchor=north] at (96,6.6) {$R_{1}$};
\end{tikzpicture}
